\documentclass[10pt]{article}
\usepackage[margin=0.6in,top=0.5in,bottom=0.5in]{geometry}
\usepackage{amsmath,amssymb}
\usepackage{enumitem}
\usepackage{booktabs}
\usepackage{tabularx}
\usepackage{xcolor}
\usepackage{hyperref}
\usepackage{titlesec}
\titlespacing*{\section}{0pt}{0.4em}{0.15em}
\titleformat{\section}{\normalfont\bfseries\normalsize}{}{0pt}{}
\setlist{nosep,leftmargin=1.2em}
\renewcommand{\arraystretch}{1.0}
\hypersetup{colorlinks=true,urlcolor=blue!50!black,linkcolor=black,citecolor=black}
\pagestyle{empty}

\begin{document}

\begin{center}
{\large\bfseries Technical Summary Sheet} \quad
6.5930 / Hardware Architectures for Deep Learning, Spring 2026\\[1pt]
\textbf{Topology-Aware Mapping for Photonic Interconnects in TPU-Class LLM Accelerators}\\
Pranav Sitaraman (\texttt{sitaraman.pranav@gmail.com}) \quad and \quad
Simone Surdi (\texttt{simonesurdi@gmail.com})
\end{center}

\section*{Repository}
\texttt{final-project/workspace/final\_project/} (extends the
AccelForge cluster example). Reproduce all numbers with
\texttt{make study \&\& make notebook \&\& make test} from
\texttt{final-project/}. The 71-test suite is the canonical
correctness check; nine figures live under
\texttt{results/presentation/}.

\section*{Three Key Insights / Innovations}
\begin{enumerate}
\item \textbf{The intra-chip mapping is stable; the distributed
schedule is what changes.} For the 64-chip TPU cube, AccelForge
converges to the same partitioning/dataflow/tile signature for all
seven technologies (\textbf{0/24} inner mappings change). The benefit
lives one layer up: \textbf{5/16} workload--scenario pairs flip their
optimal communication pattern, and the optimal compiler-visible
schedule changes whenever multicast is available. Separating the two
layers matches the 6.5930 mapping formalism rather than conflating it.

\item \textbf{Hardware multicast, not raw bandwidth, dominates
shared-state workloads.} A 3.5~Tb/s, 6~pJ/bit advanced chiplet
baseline---chosen to defeat the bandwidth hypothesis---is still
beaten by M1000 by \textbf{99.85\%} on shared-DRAM multicast: $N$
unicasts collapse to one broadcast, and contention disappears on a
switched fabric. The same chiplet is competitive on node-to-node
exchange.

\item \textbf{Photonics is not a uniform multiplier.} On a 9216-chip
Ironwood-class superpod, long-context inference reaches
\textbf{14.50}$\times$ speedup, but training only reaches
\textbf{1.43}$\times$. Once collectives shrink, training becomes
compute-dominated; inference stays network-exposed because decode and
MoE routing cross the fabric every step. The strong case for
photonics is inference and long-context serving, not training payoff.
\end{enumerate}

\section*{Headline Results (reproducible from \texttt{results/})}
\begin{tabularx}{\linewidth}{@{}lrX@{}}
\toprule
Axis & Value & Setting \\
\midrule
16-node energy reduction        & 72.21\%       & EVK100 multicast \\
16-node latency reduction       & 97.57\%       & EVK100 multicast \\
TPU-cube network energy red.    & 99.91\%       & M1000 multicast \\
TPU-cube AccelForge energy red. & 77.70\%       & M1000 KV attention \\
Inner mapping signatures        & 0/24 changed  & TPU cube \\
Distributed maps changed        & 5/16 cases    & cluster study \\
Compiler-schedule best gain     & 67.75\%       & M1000 remote streaming \\
Multicast scaling exponent      & 0.00          & M1000 vs torus / chiplet \\
Chiplet-relative reduction      & 99.85\%       & M1000 multicast \\
Superpod inference speedup      & 14.50$\times$ & Deep Think proxy \\
Superpod training speedup       & 1.43$\times$  & Pro MoE proxy \\
Superpod inference energy red.  & 51.07\%       & Deep Think proxy \\
\bottomrule
\end{tabularx}

\section*{Contributions}
\textbf{Simone Surdi} -- AccelForge/network model: topology classes
(ring, mesh, 3-D torus, optical crossbar / switched fabric); the
\texttt{LogicalClusterMemory} extension; the iterative network-aware
feedback loop (\texttt{solve\_iteratively}); mapping-signature
extraction; cluster, TPU-cube, and chiplet-baseline studies;
technology-catalog constants and sensitivity sweeps.
\quad
\textbf{Pranav Sitaraman} -- Compiler/superpod analysis: schedule
enumeration (SPMD collectives, remote HBM streaming, replicated HBM
cache, multicast cache fill, MoE locality, full photonic stack);
the 9216-chip Ironwood-class stage model (prefill / decode / refill /
MoE / training / gradient / optimizer); Gemini-class proxy workloads;
ablation ladder; integration tests; report, poster, and presentation
deliverables.

\section*{Reproducibility and Threats to Bias}
Lightmatter constants are from public product pages and SC25/ISES
talks; TPU constants are from Cloud TPU docs and the Ironwood
blog / Hot Chips tutorial. EVK100 is reported both as the conservative
3.2~pJ/bit product target and the 1.9~pJ/bit Passage-family
sensitivity point; M1000 uses 114.6~Tb/s bidirectional and 2.3~pJ/bit
including laser power. Optical multicast is modeled as one logical
broadcast (optimistic but explicit). The superpod stage model is
analytical with an overlap factor; cycle-accurate simulation is left
to future work. Gemini parameter counts are unpublished, so we use
public proxy stress workloads. The 71-test suite asserts technology
constants, the multicast scaling exponent, and the
schedule-changes-under-photonics property as regression tests.

\end{document}
